\section{概率公理与概率空间}
\label{sec:05-Probability/01-Probability-Spaces/02}

一次故障演练评审，团队给三种情况各打了一个``发生可能性''：磁盘故障 \(0.3\)，网络抖动 \(0.2\)，两者同时出问题 \(0.1\)。有人接着问，``至少出一样''该打多少？屋里报出的数从 \(0.5\) 一路散到 \(0.25\)。

先别管哪个对。留意一件更基础的事：这些数是人随手指派的，没有任何机制保证它们彼此相容。真有人给``至少出一样''打了 \(0.25\)，模型立刻自己打起来——磁盘故障单独发生就已经属于``至少出一样''，一块小区域的分数反而高过包住它的大区域。这类矛盾不需要复杂场景就能造出来，也不会有人在算式里替你发现。

上一节把三件套的前两件装好了：\(\Omega\) 说清一次试验能产出哪些结果，\(\mathcal F\) 说清哪些集合算事件。现在轮到第三件，往每个事件上贴一个数。问题是贴数字这件事自由度太大，得先约束住。Kolmogorov 在 1933 年给出的约束只有三条，其余的规矩全是从这三条推出来的定理。

\subsection{三条公理，以及为什么只要三条}

概率空间写作 \((\Omega,\mathcal F,P)\)，其中 \(P\) 是从 \(\mathcal F\) 到 \([0,1]\) 的一个函数，满足
\[
P(A)\ge0,\qquad P(\Omega)=1,\qquad
P\Big(\bigcup_{i=1}^{\infty}A_i\Big)=\sum_{i=1}^{\infty}P(A_i)\quad(A_i\text{ 两两互斥}).
\]

第一条排除负数，第二条把总量定成 \(1\)，第三条说互不重叠的事件可以把数直接加起来，而且允许加可数无穷多次。三条都对应得上日常直觉，唯一超出直觉的是第三条里的``无穷''，稍后单独交代。

评审会上那份打分之所以出问题，是因为它违反了从这三条推出来的单调性。有了公理，剩下的规则不必再靠人协商，也不会出现两条规矩打架的局面——凡是能从三条推出来的都成立，凡是推不出来的就不是概率论的承诺，得由建模者另行说明。

把 \(\Omega\) 想象成一块面积为 \(1\) 的方板，事件是板上的一块区域，\(P(A)\) 就是这块区域的面积。三条公理在这幅图里是三句大白话：面积非负；整块板的面积是 \(1\)；两块不重叠的区域，面积相加。

\begin{center}
\begin{tikzpicture}[x=1cm,y=1cm,font=\small]
  \draw[black!45,fill=black!4] (0,0) rectangle (4.6,3.2);
  \begin{scope}
    \clip (1.7,1.8) ellipse (1.2 and 1.0);
    \fill[blue!35] (3.0,1.4) ellipse (1.15 and 0.95);
  \end{scope}
  \draw[black!60] (1.7,1.8) ellipse (1.2 and 1.0);
  \draw[black!60] (3.0,1.4) ellipse (1.15 and 0.95);
  \node at (1.05,2.55) {\(A\)};
  \node at (3.75,0.75) {\(B\)};
  \node[black!60,font=\footnotesize] at (0.4,0.3) {\(\Omega\)};
  \node[black!70,font=\footnotesize] at (2.3,-0.5) {总面积 \(=1\)};

  \begin{scope}[shift={(6.6,0)}]
    \draw[black!45,fill=black!4] (0,0) rectangle (4.6,3.2);
    \fill[blue!30] (0.9,2.2) ellipse (0.75 and 0.6);
    \fill[blue!30] (2.6,2.4) ellipse (0.7 and 0.55);
    \fill[blue!30] (1.9,0.9) ellipse (0.9 and 0.6);
    \draw[black!60] (0.9,2.2) ellipse (0.75 and 0.6);
    \draw[black!60] (2.6,2.4) ellipse (0.7 and 0.55);
    \draw[black!60] (1.9,0.9) ellipse (0.9 and 0.6);
    \node[font=\footnotesize] at (0.9,2.2) {\(A_1\)};
    \node[font=\footnotesize] at (2.6,2.4) {\(A_2\)};
    \node[font=\footnotesize] at (1.9,0.9) {\(A_3\)};
    \node[black!60,font=\footnotesize] at (0.4,0.3) {\(\Omega\)};
    \node[black!70,font=\footnotesize] at (2.3,-0.5) {互不重叠时面积直接相加};
  \end{scope}
\end{tikzpicture}

\emph{概率是这块板上的面积。左图里 \(A\) 与 \(B\) 有重叠，把两块面积相加会把深色部分算两遍，扣掉一次就是容斥公式；右图里三块互不重叠，公理允许直接求和。这张图后面还会回来：讲条件概率时，``已知 \(B\) 发生''相当于把整块板裁到 \(B\) 那一块，再把它的面积重新标定为 \(1\)。}
\end{center}

\subsection{常用公式全是三条公理的推论}

先看空事件。把 \(\Omega\) 写成 \(\Omega\) 与无穷多个 \(\varnothing\) 的互斥并，可数可加性给出 \(1=1+P(\varnothing)+P(\varnothing)+\cdots\)；非负项的级数要收敛到有限值，重复出现的那一项只能是零，于是 \(P(\varnothing)=0\)。有限可加性也是这么来的：有限个互斥事件后面补上无穷多个空事件即可。

其余几条更快。把 \(\Omega\) 拆成 \(A\) 与 \(A^{c}\)，得 \(P(A^{c})=1-P(A)\)。若 \(A\subseteq B\)，把 \(B\) 拆成 \(A\) 与 \(B\setminus A\) 两块，加上非负性就有 \(P(A)\le P(B)\)；取 \(B=\Omega\) 顺带说明概率不会超过 \(1\)——这条不是第四条公理，是推论。容斥公式
\[
P(A\cup B)=P(A)+P(B)-P(A\cap B)
\]
的推导同样是两次拆分：\(A\cup B=A\cup(B\setminus A)\)，而 \(B=(A\cap B)\cup(B\setminus A)\)，消去 \(P(B\setminus A)\) 即得。上面那张面积图已经把它画完了，重叠的一块被数了两遍，扣回去一次。

这些结论没有一个出乎意料，几乎人人都能猜对。有价值的不是结论，是它们全都只用了三条起点。任何你想额外补上的``直觉规则''，要么已经在这三条里，要么和它们冲突；后一种可能被 Kolmogorov 一次性排除了。

补事件公式的实用价值值得单说一句。掷四次骰子问``至少出现一次 \(6\)''，正面枚举要分恰好一次、两次、三次、四次四种情形；换成补事件，一次都不出的概率是 \((5/6)^4\)，答案 \(1-(5/6)^4\approx0.518\)。``至少一个''几乎总该翻成``不是零个''，这是三条公理最先还给你的一笔红利。

\subsection{为什么是可数可加，不是有限可加}

有人会问，把第三条减弱成``有限个互斥事件可加''行不行？形式上行，但你会失去一大类想谈的事件。``一直掷到第一次出现 \(6\)，等待次数是某个正整数''这件事，由等待一次、两次、三次……可数无穷多个互斥事件拼成；\([0,1]\) 中``取到有理数''同样是可数多个单点的并。只有有限可加性时，这些并集的概率与各块概率之和之间没有任何关系，模型对它们无话可说。

后果比缺几个公式严重。可数可加性等价于下面这条连续性：若事件一层套一层地涨大，\(A_1\subseteq A_2\subseteq\cdots\)，并集记作 \(A\)，则
\[
\begin{aligned}
P(A)&=\sum_{i=1}^{\infty}P(B_i)\\
&=\lim_{n\to\infty}\sum_{i=1}^{n}P(B_i)
=\lim_{n\to\infty}P(A_n),
\end{aligned}
\]
其中 \(B_1=A_1\)、\(B_i=A_i\setminus A_{i-1}\) 是逐层增量，两两互斥，前 \(n\) 个的并恰好是 \(A_n\)。骨架就这一句：把``越来越大''拆成互斥的增量，再对增量用公理。若只有有限可加性，这条连续性会失效，极限号就不能穿过 \(P\)。而``当 \(n\) 很大时概率接近某个值''这类说法——大数定律、几乎必然收敛、马尔可夫链的长期行为——句句都指望它。减弱第三条，等于把整个极限定理体系的地基抽掉。

至于为什么不干脆要求``不可数可加''、以及哪些子集必须被排除在 \(\mathcal F\) 之外，属于测度论的题目，见\hyperref[ch:101]{``测度、积分与条件期望：为现代概率回填地基''}。这里只需要一个界桩：可数是能一个一个排队数完的无穷，不可数不是；公理只对前者作了承诺。

\subsection{概率为零，和不可能是两回事}

在 \([0,1]\) 上均匀取一个实数，某个特定的点比如 \(0.5\) 被取中的概率是多少？如果每个点都分到某个正数 \(\varepsilon\)，取 \(1/\varepsilon\) 个点就把总量顶破了，所以只能是零。可实验做完总会落在某一点上——一个概率为零的事件刚刚发生了。

\begin{center}
\begin{tikzpicture}[x=1cm,y=1cm,font=\footnotesize]
  \fill[blue!30] (0,0) rectangle (5.0,1.1);
  \fill[blue!16] (5.0,0) rectangle (7.5,1.1);
  \fill[blue!30] (7.5,0) rectangle (8.75,1.1);
  \fill[blue!16] (8.75,0) rectangle (9.375,1.1);
  \fill[blue!30] (9.375,0) rectangle (9.6875,1.1);
  \draw[black!55] (0,0) rectangle (10.0,1.1);
  \foreach \x in {5.0,7.5,8.75,9.375,9.6875}{\draw[black!55] (\x,0) -- (\x,1.1);}
  \node at (2.5,0.55) {\(A_1\)};
  \node at (6.25,0.55) {\(A_2\)};
  \node at (8.12,0.55) {\(A_3\)};
  \node[black!70] at (9.6,-0.45) {\(\cdots\)};
  \node[black!75] at (3.4,-0.45) {可数多块，宽度都为正，面积可以逐块加满};
\end{tikzpicture}

\emph{把区间切成可数多块，公理允许一块块把面积加起来，总和是 \(1\)。换成按``每一个点''切，块数不可数、每块宽度为零，加法失去意义——连续模型里单点概率为零，根子在这里。}
\end{center}

不可数个零加起来可以是任何数，这不是公理的漏洞，而是连续模型能工作的价格。区间的概率靠长度度量，不靠把单点概率累加；换来的好处是可以用微积分处理概率。反过来说，在 \([0,1]\) 上取到无理数的概率是 \(1\)，可取到 \(1/2\) 在逻辑上并非不可能——\(1/2\) 确实在区间里。概率论把这种情形叫几乎必然（almost surely），后面每次出现这个词，都是在提醒概率为零与不可能之间还留着一道细缝。\hyperref[sec:05-Probability/08-Limit-Theorems/01]{``随机变量的多种收敛''}里几种收敛模式的精细差别，源头正是这道缝。

\subsection{公理管什么，不管什么}

三条公理保证一个合法模型内部不自相矛盾，让常用恒等式有统一出处，也让有限与无限、离散与连续挂在同一套语言下。它们的沉默同样重要。

公理不告诉你硬币是否公平。\(P(\text{正面})\) 取 \(0.5\) 还是 \(0.37\)，公理照单全收；它管的是给定分配之后推理合不合法，不管分配本身对不对。公理也不告诉你两个事件是否独立——独立是额外加在模型上的结构，不是推论。更要紧的是，公理不替你选 \(\Omega\)。上一节那对争论下单率的同事，无论选会话还是选用户，得到的都是完全合法的概率空间；公理对哪一个更贴合业务问题不置一词。

于是最难的一步被留在了公理之外：拿到一段模糊的现实描述，怎么把它翻译成一个具体的 \((\Omega,\mathcal F,P)\)。下一节专门做这件事，重点不在翻译的结果，而在翻译途中悄悄塞进去的那几条假设。
